\section{Test Case 17: dynamic Kelvin--Voigt evaporation on the GPU}

Test Case 17 is the Kelvin--Voigt counterpart of Test Case 16.  It starts
with 100 beads and combines Yarin evaporation, nozzle insertion, collector
removal, and forced capacity growth over 1,000 integration steps.  The stored
input selects RK4; changing only \texttt{integrator} to 1 or 2 selects the
equivalent Euler or RK2 validation.

With target \texttt{nvfortran-openacc}, Euler executes one Kelvin--Voigt force
and stress evaluation, RK2 executes two, and RK4 executes four.  All
intermediate and final updates, direct Coulomb summation, evaporation, charge
smoothing and restoration, nozzle geometry, and final statistics reductions
execute on the GPU.  Kelvin--Voigt stress uses relative bead acceleration
together with the full product-rule derivative of the concentration-dependent
viscosity and modulus.  The jet state, Coulomb force, stress, and derivative
arrays remain resident between stages.

Ordinary timesteps transfer only the insertion, removal, and resize decision
scalars.  Actual topology and output events transfer their required records;
the complete active state is synchronized only for capacity reallocation and
the final checkpoint.  A runtime transfer audit confirms two full-state
rebinds for the two reallocations and no per-stage array transfers.

The accepted 1,000-step CPU and NVIDIA A30 runs for Euler, RK2, and RK4 all
produce 111 additions, 122 removals, two reallocations, and 89 final active
beads.  Three-step pre-event CPU/GPU comparisons are identical with relative
tolerance $10^{-12}$ and absolute tolerance $10^{-13}$.  Floating-point
ordering moves the first insertion threshold from step 4 on the CPU to step 5
on the GPU; the physical bending instability then amplifies this perturbation.
Consequently, equal aggregate topology, pre-event agreement, clean completion,
and the transfer audit are used instead of pointwise CPU/GPU trajectory
identity.

The historical CPU Kelvin--Voigt evaporation equation omits aerodynamic drag
and lift even when \texttt{airdrag yes} is present.  The GPU implementation
preserves this established semantics.  The development-only target
\texttt{nvfortran-openacc-kv-host-forces} evaluates every force stage through
the trusted host equations and uploads its derivatives.  This diagnostic is
useful for isolating floating-point differences but is unsuitable for
performance measurements.  The input is stored in
\texttt{examples/input-17/input.dat}.

The paired CPU/GPU checks for all three deterministic integrators are
automated by
\texttt{tests/performance/dynamic/validate\_evaporation.sh}; see Test Case 16
for the validation criteria and build policy.
